\section{Discussion}
\label{sec:discussion}
% Written from the final tables (results/final/stats_held.json, results/bdcxr/stats.json, results/final_hybrid_stats.json,
% results/final_xm_stats.json). Failed hypotheses are named first.
\paragraph{Principal findings} Several pre-specified results did not hold. The soft-tissue image alone did not improve any downstream task
(Tables~\ref{tab:downstream_cls} and~\ref{tab:downstream_det}); a multi-stream classifier that read the radiograph, soft-tissue and bone
images did not improve on the radiograph alone (H6, Table~\ref{tab:cls_primary}); SSIM fell short of the reference model by more than its
margin (Table~\ref{tab:h2}); and conformal intervals calibrated on synthetic data covered 14\,\% instead of 90\,\% of real bone pixels
(Table~\ref{tab:reliability}), 89\,\% after an exploratory recalibration on 24 real pairs (Table~\ref{tab:bone_recal}). Real-nodule contrast
was higher than with the reference model on JSRT and lower on Node21 (Table~\ref{tab:nodule_cnr}). In Study III, distillation from CT did not
improve real-radiograph classification (H11, H12), the CT and X-ray Grad-CAM++ maps did not agree (H14), and the Swin classifier trailed
DenseNet-121 over the 14 diseases (H13), almost entirely through one rare class and with a reference trained twice as long at four times
the encoder learning rate (Table~\ref{tab:xm_hyp}). The central results of Study I held.
Lesion absorption fell steadily with the twin weight (Table~\ref{tab:dose_response}); on an audit set no model had seen during development,
the twin ensemble absorbed a third of what the released CT-supervised model absorbed and less than a third of what the matched control
absorbed (Table~\ref{tab:audit_contrasts}), at non-inferior bone-region fidelity; and the effect replicated with two other encoders
(Table~\ref{tab:encoders}) and on paediatric radiographs (Section~\ref{sec:bdcxr}).

\paragraph{What the audit adds} Fidelity metrics and lesion absorption rank suppressors differently (Figure~\ref{fig:pareto}). The v1 model
has the highest bone-region PSNR of the CT-supervised models and the highest absorption; ours and the reference model are within 0.06\,dB of
each other and a factor of three apart in absorption. BS-LDM, trained on dual-energy pairs from the evaluation hospital, is about 7\,dB ahead
in fidelity on that hospital's images, yet absorbed 0.42 of an inserted lesion in the development audit (at unpaired sites, so only
descriptive): fidelity to a dual-energy target does not protect what the target never showed. Absorption must still be read beside
suppression strength. The reference model removes more rib-edge energy than ours (48\,\% against 30\,\%) and more non-rib energy (22\,\%
against 8\,\%), so part of its higher absorption may reflect stronger suppression in general.

\paragraph{What the regulariser costs} On real paired radiographs the twin term costs no measurable bone-region fidelity up to $w=16$ (21.22
to 21.28\,dB across $w=0\ldots16$); at $w=32$ and 64 PSNR drops by about 0.3\,dB, so $w=16$, chosen on development data before the freeze,
is the usable operating point. The cost appears where the truth is exact: on the Qure-projection test split the bone error of the twin
ensemble is 41\,\% higher than that of the control (MAE 0.0041 against 0.0029). The regulariser thus trades synthetic bone accuracy for
lesion invariance, a trade not visible on real images. The ablation shows that the effect needs gradients through both twin branches and a
consistency term confined to the pasted lesion; showing lesions during training without the twin term (CLIA) did not lower absorption. The
reference model is a SwinV2-B at $1024$~px trained with a far larger budget, so only the matched no-twin control, which absorbs slightly more
than the reference model (0.178 against 0.149), isolates the regulariser. Diffusion models trained on real dual-energy pairs
\citep{sun2026bsldm,sun2026debonedit} need paired data from the target site, sample slowly (about 2.5 minutes per image for BS-LDM on a T4
GPU) and were not designed for lesion invariance.

\paragraph{Do classifiers gain from suppressed inputs?} The multi-stream hybrid matched the radiograph alone on all three datasets; on NIH the
interval rules out a gain larger than 0.0024 in macro-AUROC, so the null result is informative. The hybrid also matched the capacity
control, and plain channel stacking did at least as well as separate streams. Eq.~\eqref{eq:info} offers one explanation: the suppressor adds
a transformation learned from CT but no information, and a classifier trained on thousands of radiographs may learn whatever bone-invariant
features its labels reward. Nodules and masses had the largest estimates on both adult datasets ($+0.008$ to $+0.009$ for the soft-tissue
input), and the pointing game agreed: maps of the radiograph stream found focal lesions less often than other findings, whereas maps of the
soft-tissue stream found both about equally often. Two cautions apply. The NIH test split had already been used by the downstream classifiers
of Study I, so the NIH result is a re-test with new models; and VinBigData, the independent test, has a single merged nodule/mass class
whose interval included zero. The loss for diffuse findings appeared on NIH only, mostly for infiltration; with report-derived NIH labels
\citep{wang2017chestxray} and no such loss on the radiologist-labelled VinBigData, these data cannot tell whether suppression removes cues
for diffuse disease or interacts with label noise.

\paragraph{Why CT did not teach the radiograph classifier} The distilled classifier matched the one started from ImageNet within
$\pm0.0016$ AUROC on the shared findings and matched the classifier pretrained on the same projections without the teacher. First, the
teacher had little to give: on the six findings it was hardly more accurate on CT (0.821) than P0 on the projections of the same test
patients (0.814), and it was weakest for lung lesions (0.608), where CT should help most. Second, the 2.5D recipe was fragile, with one
collapsed and one late-learning run in seven, so students learned from teachers of uneven quality. Third, gains on the projections had to
survive the move to real radiographs: the zero-shot results (0.640--0.648 on NIH) show how wide that gap is, and ten epochs of fine-tuning on
\num{77890} radiographs may override the initialisation, as the equal results of F$_\mathrm{p0}$ and F$_\mathrm{in}$ suggest. The negative
result therefore concerns this teacher and this amount of paired CT, not cross-modal distillation in general.

\paragraph{Explanations and robustness} The two modalities did not produce maps of the same kind. The teacher's Grad-CAM++ mainly ranks
slices, so its projected map varies from top to bottom but hardly across, while the $8\times8$ token grid of the last SwinV2-T stage gives
Grad-CAM++ maps that spread along the image borders. On real radiographs Swin Grad-CAM++ hit expert boxes less than half as often as
DenseNet-121 Grad-CAM++, whereas Integrated Gradients and decoder attention localised about as well for both encoders; agreement between
modalities should therefore be measured with maps that localise. Robustness, in turn, was not visible in accuracy: the Swin classifiers were
not more accurate than DenseNet-121, but noise flipped about a third as many finding-free images (0.27--0.29 against 0.78). Under rib
emphasis the distilled classifier flipped fewer images than the one started from ImageNet ($p=0.044$, uncorrected), but so did the
classifier pretrained on projections without the teacher, which points to pretraining rather than to CT.

\paragraph{Uncertainty} The bone-image intervals held where they were calibrated and failed under the shift to real radiographs; 24 real pairs
from the target site sufficed to restore coverage here, although transfer between sites is untested. For the classifiers, conformal risk
control kept its guarantee for every input and encoder, but label sets holding about half of the 14 NIH labels are of little use to a
reader, and the suppressor's own uncertainty made abstention worse in every configuration: uncertainty about the bone image is not
uncertainty about the diagnosis.

\paragraph{Recommendations and future work} A bone suppressor should be audited for lesion absorption, reported next to a suppression
measure, before it is used; suppressed images should be shown beside the original rather than instead of it; and the ensemble spread should
serve as a warning flag rather than a calibrated interval unless recalibrated on real pairs (it rose 2.7-fold on paediatric radiographs).
Classifiers should not be expected to gain more than small amounts of overall accuracy from bone-suppressed inputs, which helped only as
extra channels, or from a CT teacher of this kind, and any benefit should be tested per finding against a control of the same capacity. Swin classifiers of radiographs should be explained with Integrated
Gradients or decoder attention rather than Grad-CAM++, with maps scored against expert boxes, and claims of robustness should rest on
perturbations that add no disease. Future work should test prediction at $1024$~px, recalibration on real pairs from more sites, other
lesion types and readers, and, for transfer from CT, a stable volumetric teacher trained on many more volumes.
